%% ma: L12-B14-0075
%% lop: 12
%% chuong: 5
%% bai: 14
%% dang: 12.14.3
%% loai: TN4
%% muc: TH
%% dap_an: "C"
%% nguon: "To_n_12_PT_M_t_Ph_ng.pdf (Đề kiểm tra đánh giá kiến thức, Chuyên đề Phương trình mặt phẳng, Nguyễn Anh Khôi - Lê Minh Khôi), Câu 30"
%% kien_thuc: "Bài 14 (tích có hướng), Bài 8 (tọa độ vectơ)"
%% trang_thai: chinh-thuc
%% ghi_chu: "Tài liệu gốc không kèm đáp án; đáp án do Claude giải và kiểm bằng sympy."
\begin{ex}
Trong không gian $Oxyz$, cho ba điểm $A(1;3;2)$, $B(2;-1;5)$ và $C(3;2;-1)$. Gọi $\vec n=\left[\overrightarrow{AB},\overrightarrow{AC}\right]$ là tích có hướng của hai vectơ $\overrightarrow{AB}$ và $\overrightarrow{AC}$. Tìm tọa độ vectơ $\vec n$.
\choice{$\vec n=(3;-9;9)$}{$\vec n=(9;7;15)$}{\True $\vec n=(15;9;7)$}{$\vec n=(9;3;-9)$}
\loigiai{$\overrightarrow{AB}=(1;-4;3)$, $\overrightarrow{AC}=(2;-1;-3)$. $\vec n=\left(12+3;\,6+3;\,-1+8\right)=(15;9;7)$. Chọn C.}
\end{ex}
